
\documentclass[11pt]{article}

\usepackage[margin=1in]{geometry}
\usepackage{booktabs}
\usepackage{graphicx}
\usepackage{array}
\usepackage{hyperref}
\usepackage{enumitem}

\title{Restaurant Vibe Classification with Locally Fine-Tuned DistilBERT}
\author{Natalie Huante}
\date{May 2026}

\begin{document}

\maketitle

\begin{abstract}
% Briefly summarize:
% - the restaurant vibe classification problem
% - Yelp Open Dataset usage
% - DistilBERT fine-tuning
% - baseline comparison
% - final results and conclusions
\end{abstract}

\section{Introduction}

% Suggested structure:
% Paragraph 1:
% - restaurant discovery challenges
% - subjective restaurant atmosphere / “vibe”
% - importance of reviews and user experiences
%
% Paragraph 2:
% - NLP for understanding restaurant reviews
% - transformer-based text classification
%
% Paragraph 3:
% - project overview
% - local fine-tuning goals
%
% Contributions:
% - custom restaurant vibe dataset
% - local DistilBERT fine-tuning
% - baseline comparison
% - semantic error analysis

\section{Problem Definition}

% Define the NLP task clearly:
% Input:
% - restaurant categories
% - Yelp review text
%
% Output:
% - one of seven vibe labels
%
% Mention:
% - supervised multi-class text classification

\subsection{Vibe Labels}

\begin{table}[h]
\centering
\begin{tabular}{ll}
\toprule
Label & Description \\
\midrule
study\_spot & Quiet cafes and work-friendly spaces \\
late\_night & Bars, nightlife, and evening atmosphere \\
casual\_hangout & Relaxed and informal dining \\
family\_friendly & Restaurants suited for families and groups \\
upscale & Elegant or fine dining experiences \\
trendy & Modern, stylish, or social restaurants \\
date\_night & Romantic or intimate dining atmosphere \\
\bottomrule
\end{tabular}
\caption{Restaurant vibe classification labels}
\end{table}

\section{Dataset}

\subsection{Data Source}

% Describe Yelp Open Dataset:
% - business.json
% - review.json
% - metadata fields
% - review text

\subsection{Dataset Filtering}

% Describe:
% - restaurant category filtering
% - removal of irrelevant businesses
% - geographic filtering
% - Santa Barbara / Goleta / Carpinteria focus

\subsection{Dataset Construction}

% Explain:
% - merging business metadata with reviews
% - constructing input_text
%
% Example format:
% Categories: Coffee & Tea, Cafes | Review: Quiet cafe with students studying.

\subsection{Labeling Process}

% Discuss:
% - initial imbalance problem
% - keyword candidate pools
% - manual balancing strategy
% - final 200-example dataset

\subsection{Train / Validation / Test Split}

% Include:
% - 70/15/15 split
% - stratified splitting
% - balanced label distribution

\section{Model and Fine-Tuning Method}

\subsection{Model Selection}

% Explain:
% - DistilBERT selection
% - local hardware feasibility
% - compact transformer architecture
% - computational efficiency

\subsection{Tokenization}

% Mention:
% - DistilBertTokenizer
% - padding
% - truncation
% - max sequence length

\subsection{Fine-Tuning Method}

% Explain:
% - sequence classification head
% - Hugging Face Trainer
% - PyTorch
% - local training
%
% IMPORTANT:
% Mention that all trainable model parameters were updated.

\subsection{Baseline Model}

% Describe:
% - TF-IDF vectorization
% - Logistic Regression classifier
% - purpose of baseline comparison

\section{Experimental Setup}

% Include:
% - local hardware setup
% - software libraries
% - PyTorch
% - Hugging Face Transformers
% - scikit-learn

\subsection{Training Configurations}

\begin{table}[h]
\centering
\begin{tabular}{llll}
\toprule
Experiment & Learning Rate & Epochs & Notes \\
\midrule
Model 1 & 3e-5 & 8 & Baseline fine-tuning \\
Model 2 & 2e-5 & 8 & Warmup scheduling \\
Model 3 & 3e-5 & 10 & Extended training duration \\
\bottomrule
\end{tabular}
\caption{DistilBERT experiment configurations}
\end{table}

\section{Results}

% Discuss:
% - overall model performance
% - best model selection
% - baseline comparison
% - generalization performance

\subsection{Final Model Comparison}

\begin{table}[h]
\centering
\begin{tabular}{lcc}
\toprule
Model & Accuracy & Weighted F1 \\
\midrule
TF-IDF + Logistic Regression & 0.67 & 0.66 \\
DistilBERT (Model 1) & 0.67 & 0.68 \\
\bottomrule
\end{tabular}
\caption{Baseline vs. fine-tuned DistilBERT performance}
\end{table}

\subsection{Confusion Matrix}

% Insert confusion matrix figure here

\section{Error Analysis}

% Discuss:
% - semantically overlapping classes
% - trendy vs date_night
% - upscale vs trendy
% - late_night vs trendy
%
% Mention:
% - many errors were semantically reasonable
% - contextual ambiguity

\subsection{Qualitative Examples}

\begin{table}[h]
\centering
\begin{tabular}{p{4cm} p{2cm} p{2cm}}
\toprule
Input Summary & True Label & Predicted Label \\
\midrule
Quiet cafe with students studying & study\_spot & study\_spot \\
Stylish cocktail lounge & trendy & date\_night \\
Upscale wine bar & upscale & trendy \\
Family breakfast diner & family\_friendly & family\_friendly \\
\bottomrule
\end{tabular}
\caption{Example predictions from the final model}
\end{table}

\section{Ethical and Practical Considerations}

% Discuss:
% - subjectivity of vibe labels
% - labeling bias
% - geographic bias
% - Yelp review bias
% - limitations of small datasets

\section{Conclusion}

% Summarize:
% - project goal
% - local DistilBERT fine-tuning
% - baseline comparison
% - strongest findings
%
% Mention future work:
% - larger datasets
% - PEFT methods
% - recommendation systems
% - larger transformer architectures

\end{document}
